\chapter{銀河相互作用を調べる}
\label{chap:experiments}
第5章では、二つの移動するPlummerポテンシャルとテスト粒子円盤を組み合わせ、銀河相互作用を計算しました。本章では、完成したコードを数値実験の道具として使います。円盤の回転方向、近点距離、質量比を別々に変え、何が変わり、何が変わらないかを調べます。新しい重力積分器を作るのではなく、前章のKDK leapfrogを再利用します。

\begin{goalbox}[この章で学ぶこと]
\begin{itemize}
\item 変更するパラメータ、固定する条件、測定する量を、計算前に決められます。
\item 同じ初期標本を使い、prograde・retrograde、近点距離、質量比の比較ができます。
\item 近点からの経過時間をそろえ、粒子割合と形態の違いを説明できます。
\item 時間刻み、乱数標本、測定条件への依存を、物理モデルの変更と区別できます。
\end{itemize}
\end{goalbox}

\section{計算を数値実験として組み立てる}
\label{sec:ch6-design}
\subsection{まず問いを一つ決める}
「銀河を接近させると何かが起きる」という確認から、もう一歩進みます。例えば、「円盤2の回転を逆向きにすると、出身中心から遠くへ移る星の割合は変わるでしょうか」という問いを立てます。この問いに答えるには、回転以外の条件をそろえる必要があります。

同時に、何をもって違いを判断するかも決めます。本章では、同じ範囲で描いた粒子分布と、第5章の外側粒子割合を使います。図を見てから都合のよい時刻や閾値だけを選ぶのではなく、標準の測定条件を先に定めます。別の条件でも確認するときには、それを追加の試験として区別します。

\subsection{一つの変更が初期値に及ぼす影響を確認する}
独立に指定するパラメータを一つ変えても、すべての配列要素が一つだけ変わるわけではありません。近点距離を変えれば、それを実現する初期接線速度と動径速度も変わります。銀河の質量を変えれば、中心軌道だけでなく、その銀河の円運動速度も変わります。

ここでいう「一つずつ変える」は、こうした連動を止めることではありません。変更する入力と、モデルの式から計算し直す値を区別することです。表\ref{tab:ch6-design}に、本章で行う三つの系列を示します。
\begin{table}[htbp]\centering\small
\caption{基本編の比較系列。すべて第5章と同じ重力モデルを使います。質量比系列では全質量を固定しません。}\label{tab:ch6-design}
\begin{tabularx}{\linewidth}{lXX}\toprule
系列 & 直接変更するもの & 連動して作り直すもの \\\midrule
回転方向 & 円盤2の内部回転 $s_{{\rm rot},2}=+1,-1$ & 円盤2の中心に対する速度 \\
近点距離 & $D_{\rm p}=5,8,12\unit{kpc}$ & 中心の初期相対速度 \\
質量比 & $M_1$を固定した $q=M_2/M_1$ & 中心の初期位置・速度、円盤2の円運動速度 \\\bottomrule
\end{tabularx}
\end{table}

\begin{failbox}{名前だけを変えても、軌道は変わりません}
計算済みの初期状態を保ったまま、変数\texttt{d\_peri}だけを書き換えても、実際の近点距離は変わりません。近点距離は積分器へ毎ステップ与える命令ではなく、初期速度を決めるための入力です。質量を変更した場合も、必要な初期値を同じモデルから作り直します。
\end{failbox}

\section{第5章の基準と測定方法を引き継ぐ}
\label{sec:ch6-baseline}
\subsection{基準モデルを再現する}
基準モデルは、$M_1=M_2=10^{11}\Msun$、$a_1=a_2=b=3\unit{kpc}$、各円盤$N_k=1000$粒子、$R_{\rm d}=3\unit{kpc}$、$R_{\max}=15\unit{kpc}$です。中心間の初期距離は$60\unit{kpc}$、目標近点距離は$8\unit{kpc}$で、両円盤をprogradeとします。円盤生成には前章と同じseed、乱数の呼び出し順序、CDF表を使います。

内部単位はkpc、Gyr、$\Msun$です。時間積分は$0.60\unit{Gyr}$まで、刻みは$0.05\unit{Myr}$です。全粒子は2 Myrごと、中心は毎ステップ保存します。KDKの新しい加速度評価は、1計算につき$12000+1=12001$回です。各ステップで全中心・全粒子を同じ時刻へ進めます。
\inputcode{chap06/base-config}

\texttt{dict}は、パラメータを名前と値の組でまとめるPythonの辞書です。質量の配列は\texttt{base\_config}からキー\texttt{"masses"}を指定して取り出します。新しいクラスは作りません。比較用の辞書は\texttt{dict(...)}で作り、変更しない配列はそのまま参照します。本章の関数は入力配列を上書きしないため、この使い方ができます。

コードは第5章の関数定義をNotebook内にも収録しています。別の章のNotebookを先に実行する必要はありません。一方、本文では円盤のCDF、Plummer加速度、KDKの導出を繰り返しません。式と実装の確認には前章までを参照します。（第\ref{chap:interaction}章参照）

\subsection{出身中心から外へ移った割合を使う}
粒子$i$の出身銀河を$k(i)$とし、出身中心からの距離を
\begin{equation}
 s_i(t)=|\rvec_i(t)-\Rvec_{k(i)}(t)|
 \label{eq:ch6-origin-radius}
\end{equation}
とします。初期に付けた出身番号は、他方の中心へ近づいても変更しません。閾値を$s_{\rm cut}=30\unit{kpc}$とし、各出身群について
\begin{equation}
 f_k(t)=\frac{1}{N_k}\sum_{i:k(i)=k}\mathbf{1}\bigl(s_i(t)>s_{\rm cut}\bigr)
 \label{eq:ch6-fraction}
\end{equation}
を測ります。$\mathbf{1}$は条件が真なら1、偽なら0です。前章の\texttt{origin\_radii}と\texttt{outer\_fractions}を、そのまま使います。

この量は位置だけで定義した割合です。非束縛率、永久に剥ぎ取られた割合、質量損失率とは呼びません。$N_k$は重力源の質量を分割した数でもありません。例えば質量比1:10でも各円盤1000粒子を置き、両出身群を同じ標本数で表示できます。

\subsection{近点からの経過時間をそろえる}
近点距離や全質量を変えると、近点通過の時刻も変わります。同じ$t=0.40\unit{Gyr}$の図でも、一方は接近直後、他方はまだ接近中かもしれません。そこで、本章の主な比較には
\begin{equation}
 \tau=t-t_{\rm p}
 \label{eq:ch6-aligned-time}
\end{equation}
を使います。$t_{\rm p}$は、毎ステップ保存した中心間距離の最小値に対応する時刻です。連続時間の厳密な近点時刻ではありません。

比較表では$\tau=0.20\unit{Gyr}$を目標にします。星の履歴は2 Myrごとなので、目標に最も近い保存時刻を使い、実際の$t$と$\tau$も記録します。保存間隔が一定なら、この選択による目標時刻との差は1 Myr以下です。この上限は、中心の近点時刻自体の数値誤差を含む上限ではありません。
\inputcode{chap06/aligned-index}

\texttt{np.argmin}が返すのは配列内の最小値の場所です。距離の最小が計算の端にある場合や、比較したい時刻が保存範囲を外れる場合には、そのまま比較を続けず停止します。別のモデルで$\tau=0.20\unit{Gyr}$まで到達していなければ、計算時間を延ばすか、全系列の比較時刻を変更します。


% BEGIN TikZ addition: ch06_time_alignment
図\ref{fig:concept-ch6-time}で、二つの実験の近点時刻をそろえる操作を確認してください。目標の$\tau$と実際に使った保存時刻は区別し、両者を記録します。
\begin{figure}[!htbp]
\centering
% Schematic timing only. The benchmark's pericenter is measured from sampled center positions.
\begin{tikzpicture}[textbook diagram]
  \node[anchor=west,font=\sffamily\bfseries] at (0,4.73) {Compare the same stage of the encounter};
  \node[anchor=west] at (0,4.04) {(a) Time from the start};
  \foreach \y/\tag/\p/\col in {2.98/A/3.4/figblue,1.82/B/5.1/figred} {
    \node[anchor=east] at (1.0,\y) {Run \tag};
    \draw[vector] (1.25,\y) -- (11.35,\y) node[right] {$t$};
    \draw (1.45,{\y-.10}) -- (1.45,{\y+.10});\node[below,small label] at (1.45,{\y-.10}) {$0$};
    \fill[\col] (\p,\y) circle (2.2pt);
    \node[above] at (\p,{\y+.10}) {$t_{{\rm p},\tag}$};
  }
  \draw[guide] (7.95,1.5) -- (7.95,3.43);
  \node[small label,anchor=west] at (8.13,3.36) {Same $t$};
  \node[small label,text=figgray,anchor=west] at (3.10,.99) {Different times since pericenter};
  \draw[black!25] (0,.43) -- (11.80,.43);
  \node[anchor=west] at (0,-.18) {(b) Shift each pericenter to $\tau=0$};
  \foreach \y/\tag/\col/\sel in {-1.26/A/figblue/8.07,-2.50/B/figred/7.93} {
    \node[anchor=east] at (1.0,\y) {Run \tag};
    \draw[vector] (1.25,\y) -- (11.35,\y) node[right] {$\tau$};
    \fill[\col] (3.7,\y) circle (2.2pt);
    \node[below,small label] at (3.7,{\y-.10}) {$0$};
    \foreach \delta in {-.75,-.50,-.25,0,.25,.50} {\draw[black!45] ({\sel+\delta},{\y-.09}) -- ({\sel+\delta},{\y+.09});}
    \draw[\col,fill=white,line width=.9pt] ({\sel-.07},{\y-.07}) rectangle ({\sel+.07},{\y+.07});
  }
  \draw[guide] (3.7,-.67) -- (3.7,-2.78);
  \draw[guide] (7.96,-.67) -- (7.96,-2.78);
  \node[small label,above] at (7.96,-.68) {Target $\tau$};
  \node[small label,anchor=west] at (0,-3.42) {$\tau=t-t_{\rm p}$; squares: nearest stored snapshots.};
  \node[small label,text=figgray,anchor=west] at (0,-3.99) {Record the actual $t$ and $\tau$ used for each run.};
\end{tikzpicture}

\caption{開始からの時刻$t$と、近点からの経過時間$\tau=t-t_{\rm p}$の違い。下段では各近点を$\tau=0$へそろえています。小目盛りは保存時刻、四角は目標に最も近い保存時刻です。配置と目盛りは説明用で、実測の時間間隔を表していません。}
\label{fig:concept-ch6-time}
\end{figure}
\FloatBarrier
% END TikZ addition: ch06_time_alignment

\subsection{一回の実験を短い関数にまとめる}
初期値、KDK、診断という既習の順序を関数へまとめます。返す辞書には、設定と結果を一緒に入れます。比較後に、どの曲線がどの条件だったか分からなくなることを避けられます。
\inputcode{chap06/run-case}

\texttt{make\_initial\_state(**config)}は、辞書内の名前付きパラメータを関数へ渡す記法です。\texttt{*args}は、ここでは質量、Plummerスケール、有効長さの三つを順に渡します。どちらも新しい数値計算法ではなく、引数の整理です。関数が返す位置や速度を、後から「保存則に合うように」修正する操作は行いません。

\section{円盤の回転方向を変える}
\label{sec:ch6-spin}
\subsection{銀河中心の運動を固定する}
回転方向の定義は図\ref{fig:concept-ch5-spin}を参照してください。ここでも変えるのは円盤2の内部回転だけです。

基準のprograde--progradeをPP、円盤2だけをretrogradeにしたものをPRと呼びます。軌道の比角運動量は正で、反時計回りです。円盤2の回転を変えるときには、局所的な内部速度だけを反転します。
\begin{equation}
 \vvec_i'=\Vvec_2-(\vvec_i-\Vvec_2)=2\Vvec_2-\vvec_i
 \quad\bigl(k(i)=2\bigr).
 \label{eq:ch6-reverse-spin}
\end{equation}
中心の速度$\Vvec_2$は変えません。全速度に単に$-1$を掛けると、銀河全体の並進運動まで反転してしまいます。
\inputcode{chap06/reverse-spin}

同じseed・粒子数から円盤を生成すると、PPとPRの初期位置は同じで、円盤2の内部速度だけが反転します。配列から式\eqref{eq:ch6-reverse-spin}を直接計算した結果とも照合します。

本モデルでは、円盤粒子から中心への反作用がありません。また、円盤2の星から円盤1の星への力もありません。このため、PPとPRで中心軌道と円盤1の全粒子軌道は同じになります。配布コードでも配列の一致を確認します。これは一般の自己重力銀河でも成り立つ性質ではなく、本書の近似に由来します。（第\ref{chap:interaction}章参照）

\subsection{形と割合を合わせて読む}
図\ref{fig:ch6-pp}と図\ref{fig:ch6-pr}は、$\tau\simeq0.20\unit{Gyr}$での比較です。位置の表示範囲、出身色、マーカー、粒子数をそろえています。PRの円盤2では、PPと比べて遠方まで伸びた粒子の数が少なく見えます。
\begin{figure}[htbp]\centering
\includegraphics[width=0.63\linewidth]{\figpath{ch06_pp}}
\caption{PPの粒子分布。青丸と朱色三角はそれぞれ銀河1・2の出身粒子です。黒い$\times,+$はPlummer中心です。実際の比較時刻は$t=\chsixvalBaseTime\unit{Gyr}$、$\tau=\chsixvalBaseTau\unit{Gyr}$です。}\label{fig:ch6-pp}
\end{figure}
\begin{figure}[htbp]\centering
\includegraphics[width=0.63\linewidth]{\figpath{ch06_pr}}
\caption{PRの粒子分布。図\ref{fig:ch6-pp}と同じ時刻・同じ軸範囲です。円盤2の内部回転だけを変えています。円盤1と中心軌道はPPと同じです。}\label{fig:ch6-pr}
\end{figure}

定量的には、円盤2の外側割合はPPで$\chsixvalBaseTwo$、PRで$\chsixvalRetroTwo$となり、差は$\chsixvalSpinDifference$でした。これは各1000粒子の一つの標本で得た結果です。第5章の$t=0.60\unit{Gyr}$における$\chsixvalFinalPPTwo$と$\chsixvalFinalPRTwo$とは測定時刻が違うので、値を混同しません。
\begin{figure}[htbp]\centering
\includegraphics[width=0.76\linewidth]{\figpath{ch06_spin_fraction}}
\caption{円盤2の外側割合の時間変化。横軸は近点からの経過時間です。PPを青実線、PRを朱色破線で示します。粒子が出身中心から30 kpcを越えたかを数えており、非束縛判定ではありません。}\label{fig:ch6-spin}
\end{figure}

\subsection{回転方向によって相互作用が変わる理由を考える}
潮汐変形に関係するのは、星と出身中心が受ける加速度の差です。\cite{ryden2025} その差は時刻だけでなく、星が円盤のどこを動いているかにも依存します。したがって、中心軌道と初期の粒子位置が同じでも、内部速度の向きを変えれば、遭遇中に各星が受ける力の履歴は変わります。

順行では、軌道運動と同じ向きに回る星が相手の重力から受ける影響を、長く受け続ける場合があります。逆行では相手との相対的な角運動が違うため、力の向きが変わる速さも異なります。ここでは、これを計算結果の定性的な読み方とします。「すべての順行遭遇で、必ず同じ割合だけテイルが強くなる」という法則を、この二例から導くことはできません。厳密な共鳴条件の導出は本章の範囲外です。また、ここでは円盤の傾斜角を固定して共面に限定しています。回転符号の比較を、任意の三次元的な向きの比較へ一般化しません。

\section{近点距離を変える}
\label{sec:ch6-pericenter}
等質量、同じ円盤、両円盤progradeを保ち、$D_{\rm p}=5,8,12\unit{kpc}$を比べます。初期距離$D_0=60\unit{kpc}$と中心相対エネルギー$\varepsilon=0$は共通です。第5章の式から、各近点に合う中心初期速度を計算し直します。
\inputcode{chap06/pericenter-series}

$D_{\rm p}$を小さくすれば、最接近時の重力の差が大きくなり得ます。しかし、近点での速度や星の応答時間も変わります。さらに、本書の最接近では円盤の大きさが中心間距離に比べて十分小さいとは限りません。遠方での見積もり$|\delta\avec|\sim GM_2|\boldsymbol s|/D^3$を、円盤全体の厳密な応答式には使いません。（第\ref{chap:interaction}章参照）
\begin{figure}[htbp]\centering
\includegraphics[width=0.76\linewidth]{\figpath{ch06_pericenter_fraction}}
\caption{近点距離を変えた円盤2の外側割合。各曲線の近点時刻を横軸の0へそろえています。円盤・質量・回転方向・seedは共通です。}\label{fig:ch6-pericenter}
\end{figure}

今回の$\tau\simeq0.20\unit{Gyr}$では、円盤2の割合は近点5、8、12 kpcの順に$\chsixvalPeriSeries$でした。\textbf{この指標は、近点を小さくするほど単調に大きくなったわけではありません。} したがって、「より近いほど、30 kpcより外の粒子割合が必ず増える」とは結論できません。

この非単調性を、ただちに数値計算の失敗とすることもできません。測っているのは特定時刻・特定閾値での粒子数です。星の位相、遭遇の速さ、その後の運動などをまとめて反映します。どの効果が主な原因かを分けるには、さらに軌道やトルクの履歴を調べる必要があります。本章では、まず同じ条件で得た結果を記録し、予想と違うことを消さない姿勢を優先します。

\subsection{中心間の有効長さも確認する}
近点距離を比べる際には、中心間力の有効長さ$b$も無視できません。同じ$D$と相手の質量で、中心が受ける有効加速度の大きさ$A_{\rm eff}$と点質量の加速度$A_{\rm point}$を比べると、
\begin{equation}
\frac{A_{\rm eff}}{A_{\rm point}}=\frac{GM D/(D^2+b^2)^{3/2}}{GM/D^2}=\frac{D^3}{(D^2+b^2)^{3/2}}
\label{eq:ch6-softening-ratio}
\end{equation}
です。表\ref{tab:ch6-softening}は、既存の系列で使った$b=3\unit{kpc}$に対してこの比を計算したものです。$D=5\unit{kpc}$では比は約$\chsixvalSoftFive$で、点質量の値より小さくなっています。
\begin{table}[htbp]\centering\small
\caption{同じ距離での中心間加速度の比較。軌道の追加計算ではなく、式\eqref{eq:ch6-softening-ratio}の評価です。$b=3\unit{kpc}$を固定しています。}\label{tab:ch6-softening}
\begin{tabular}{rrr}\toprule
$D$ [kpc] & $\sqrt{D^2+b^2}$ [kpc] & $A_{\rm eff}/A_{\rm point}$ \\
\midrule
5 & 5.831 & 0.631 \\
8 & 8.544 & 0.821 \\
12 & 12.369 & 0.913 \\
\bottomrule\end{tabular}

\end{table}

ただし、これは中心の\textbf{並進加速度}の比です。星と中心の加速度差を表す潮汐力の比でも、遭遇を通じた力積やテイルの長さでもありません。$b$は中心軌道を、$a_1,a_2$は星が受ける場を決めます。これらの有限長さは結果に関係する候補ですが、表だけから外側割合の非単調性の原因を特定することはできません。

本章では近点系列を変更せず、この力の比較を演習6.9にします。$b$を変えることは、時間刻みを細かくする操作ではなく、中心軌道の物理モデルを変える操作です。標準の$b=a_1=a_2$から$b$だけを変えれば、中心と星が同じ相手場を感じる対応も変わります。原因を切り分ける追加実験では、その変更まで記録する必要があります。

\section{質量比を変える}
\label{sec:ch6-mass}
\subsection{主銀河の質量を固定する}
$q=M_2/M_1$とし、$M_1=10^{11}\Msun$を固定して$q=1,0.5,0.1$を比べます。第2章の「全質量を固定して質量比を変える実験」とは異なります。今回は全質量が変わるため、相対軌道に入る$G(M_1+M_2)$も変わります。

円盤2の円運動速度は$\sqrt{M_2}$に比例して変わります。初期相対速度を決めるときにも、新しい全質量を使います。一方、Plummerスケール長、指数円盤の大きさ、各粒子数は同じです。これは、一定のサイズで質量を変更した教材用の系列です。実際の銀河が従う質量とサイズの関係を再現した系列ではありません。
\inputcode{chap06/mass-series}

\subsection{近点通過前後を同じ段階で比べる}
表\ref{tab:ch6-timing}に、全系列で実際に使った時刻を示します。$q=0.1$では近点が約$\chsixvalQsmallTime\unit{Gyr}$へ遅れています。$t=0.47\unit{Gyr}$だけを比較すれば、基準モデルより近点後の経過時間が短い状態を見てしまいます。
\begin{table}[htbp]\centering\small
\caption{比較系列と測定時刻。P5・P12は近点距離、Q0.5・Q0.1は質量比を変更したものです。$t_{\rm used}$は星の最寄り保存時刻です。}\label{tab:ch6-timing}
\begin{tabular}{lrrrrr}\toprule
ケース & $q$ & $D_{\rm p}$ [kpc] & $t_{\rm p}$ [Gyr] & $t_{\rm used}$ [Gyr] & $\tau_{\rm used}$ [Gyr] \\
\midrule
PP & 1.0 & 8 & 0.26990 & 0.470 & 0.20010 \\
PR & 1.0 & 8 & 0.26990 & 0.470 & 0.20010 \\
P5 & 1.0 & 5 & 0.25470 & 0.454 & 0.19930 \\
P12 & 1.0 & 12 & 0.28750 & 0.488 & 0.20050 \\
Q0.5 & 0.5 & 8 & 0.31165 & 0.512 & 0.20035 \\
Q0.1 & 0.1 & 8 & 0.36390 & 0.564 & 0.20010 \\
\bottomrule\end{tabular}

\end{table}
\begin{figure}[htbp]\centering
\includegraphics[width=0.76\linewidth]{\figpath{ch06_mass_fraction}}
\caption{質量比系列の円盤2の外側割合。主銀河の質量は固定しています。粒子数は各円盤1000のままで、質量を粒子数で表していません。}\label{fig:ch6-mass}
\end{figure}

\subsection{一つの傾向で全結果を説明しない}
表\ref{tab:ch6-response}では、主銀河側の$f_1$は、この三つの質量比で$\chsixvalMassSeriesOne$でした。一方、円盤2の$f_2$は$\chsixvalMassSeriesTwo$で、やはり単調ではありません。小質量の銀河は自分の場による束縛が弱くなりますが、中心軌道や内部回転も同時に変わっています。この系列だけで、質量比と外側割合の普遍的な関係を決めないようにします。

$f_1=0$は、この標本・時刻で30 kpcを越えた星がなかったことです。円盤1の星が全く摂動を受けなかったことや、円盤内部に何の変形もなかったことを意味しません。閾値で集約した量が捨てている情報も意識します。
\begin{table}[htbp]\centering\small
\caption{外側粒子割合の比較。左2列は近点後の段階をそろえた値、右2列は開始から同じ0.60 Gyrの値です。同じ行でも、異なる時刻の測定です。}\label{tab:ch6-response}
\begin{tabular}{lrrrr}\toprule
ケース & $f_1(\tau\simeq0.20)$ & $f_2(\tau\simeq0.20)$ & $f_1(t=0.60)$ & $f_2(t=0.60)$ \\
\midrule
PP & 0.390 & 0.401 & 0.425 & 0.432 \\
PR & 0.390 & 0.055 & 0.425 & 0.065 \\
P5 & 0.331 & 0.352 & 0.381 & 0.388 \\
P12 & 0.304 & 0.299 & 0.330 & 0.322 \\
Q0.5 & 0.225 & 0.544 & 0.272 & 0.584 \\
Q0.1 & 0.000 & 0.427 & 0.002 & 0.508 \\
\bottomrule\end{tabular}

\end{table}

\section{差が何に由来するかを確かめる}
\label{sec:ch6-validation}
\subsection{測る閾値を変えてみる}
同じ計算結果を使い、$s_{\rm cut}=20,30,40\unit{kpc}$でPPとPRを比べます。これは運動方程式の再計算ではなく、測定条件の変更です。時刻と粒子の軌道は同じに保ちます。
\inputcode{chap06/threshold-check}
\begin{table}[htbp]\centering\small
\caption{閾値を変えたPP--PR比較。近点後約0.20 Gyr、円盤2の出身粒子についての測定です。}\label{tab:ch6-threshold}
\begin{tabular}{rrrr}\toprule
$s_{\rm cut}$ [kpc] & $f_{2,\rm PP}$ & $f_{2,\rm PR}$ & PP $-$ PR \\
\midrule
20 & 0.447 & 0.092 & 0.355 \\
30 & 0.401 & 0.055 & 0.346 \\
40 & 0.314 & 0.038 & 0.276 \\
\bottomrule\end{tabular}

\end{table}

今回の三つの閾値では、どれもPPの方が大きいという傾向は同じでした。しかし、差の大きさは変わります。報告するときは「PPの方が何\%大きい」とだけ書かず、出身中心からの距離、閾値、時刻、分母を明記します。

\subsection{乱数標本を対にして比べる}
次はseedを42、43、44へ変えます。各seedについて、PPとPRに同じ初期粒子位置を使います。二つの計算で別々の乱数を引くよりも、回転方向による差と標本の違いを分けやすくなります。
\begin{table}[htbp]\centering\small
\caption{三つのseedでのPP--PR比較。各行が同じ標本から作った一組です。全行で$\tau\simeq0.20\unit{Gyr}$、$s_{\rm cut}=30\unit{kpc}$です。}\label{tab:ch6-seed}
\begin{tabular}{rrrr}\toprule
seed & $f_{2,\rm PP}$ & $f_{2,\rm PR}$ & PP $-$ PR \\
\midrule
42 & 0.401 & 0.055 & 0.346 \\
43 & 0.409 & 0.060 & 0.349 \\
44 & 0.415 & 0.054 & 0.361 \\
\bottomrule\end{tabular}

\end{table}

この三組では、回転方向による差はseedを変えても残りました。ただし、三つのseedだけで統計分布や信頼区間を精密に推定したとは言いません。また、初期の有限標本を変える試験は、数値積分の刻みを変える試験でも、現実の銀河の多様性をすべて調べる試験でもありません。

\subsection{新しい条件でも時間刻みを確認する}
第5章で標準条件を検証していても、別の近点や質量比で同じ誤差になるとは限りません。本章ではPP、近点5 kpc、質量比0.1について、$\Delta t=0.05$と$0.025\unit{Myr}$を比べます。初期配列は作り直さず、同じ配列を渡します。全301個の保存時刻も一致させます。

二つの数値解の位置差を、中心を除いた全粒子と全保存時刻について調べます。
\begin{align}
 \delta_{\max}&=\frac{\max_{j,i}|\rvec_{i,h}(t_j)-\rvec_{i,h/2}(t_j)|}{R_{\rm d}},\\
 \delta_{\rm rms}&=\frac{1}{R_{\rm d}}\sqrt{\frac{1}{N_{\rm save}N}\sum_{j,i}|\rvec_{i,h}(t_j)-\rvec_{i,h/2}(t_j)|^2}.
 \label{eq:ch6-position-differences}
\end{align}
$h$は粗い側の刻みです。これらは真の解に対する誤差ではなく、二つの数値解の差です。$\Delta f_k$は同じ比較時刻での細かい計算の$f_k$から、粗い計算の$f_k$を引いた量とします。近点の推定時刻を細かい側で選び直して比較時刻まで動かすことはせず、粗い側で選んだ$t_{\rm used}$にそろえます。
\begin{table}[htbp]\centering\small
\caption{0.05/0.025 Myrの刻み比較。位置差は全301時刻、$\Delta f_k$は表\ref{tab:ch6-timing}の同一時刻で測定します。$e_{E,\rm c}$は第5章と同じ尺度で規格化した中心軌道のエネルギー誤差です。}\label{tab:ch6-resolution}
\begin{tabular}{lrrrrr}\toprule
ケース & $\delta_{\max}$ & $\delta_{\rm rms}$ & $\Delta f_1$ & $\Delta f_2$ & 細かい側の $e_{E,\rm c}$ \\
\midrule
PP & $1.61\times 10^{-2}$ & $1.67\times 10^{-4}$ & +0.000 & +0.000 & $8.26\times 10^{-7}$ \\
P5 & $1.25\times 10^{-2}$ & $2.07\times 10^{-4}$ & +0.000 & +0.000 & $3.98\times 10^{-6}$ \\
Q0.1 & $5.14\times 10^{-4}$ & $1.19\times 10^{-5}$ & +0.000 & +0.000 & $4.54\times 10^{-7}$ \\
\bottomrule\end{tabular}

\end{table}

今回の$\Delta f_k$はすべて0でした。しかし、位置差は0ではありません。最大位置差と二乗平均平方根にも大きな違いがあります。全体の割合だけなら一致していても、一部の星の軌道は刻みに敏感です。外側割合を議論するのか、一つの星の詳細な軌道を議論するのかによって、必要な確認の粒度が変わります。

この二つの刻みだけから収束次数を測ったとは言いません。第5章の複数刻みでの試験と、本章の変更条件での比較を組み合わせています。さらに小さな近点や長い計算を追加する場合は、必要に応じてもう一段細かい刻みでも確認してください。

\begin{failbox}{指標の一致を、すべての軌道の一致としない}
$f_k$は閾値を越えたかどうかだけを数えます。同じ側にとどまる星の位置が変わっても値は変わりません。逆に、閾値に近い1粒子が反対側へ移るだけで$1/N_k$変わります。集約した指標、個別位置の差、中心の保存量を、それぞれの意味で確認します。
\end{failbox}

\section{小さな研究としてまとめる}
\label{sec:ch6-report}
前半の最終課題は、三つの系列から一つを選び、基準と変更例の差を短いレポートへまとめることです。多くのパラメータを走査する必要はありません。まず一つの問いに対して、条件をそろえた比較ができることを目指します。
\begin{table}[htbp]\centering\small
\caption{最終レポートに必要な内容。図を作るだけでなく、どのモデルの何を測った結果かを記録します。}\label{tab:ch6-report}
\begin{tabularx}{\linewidth}{lX}\toprule
項目 & 記述する内容 \\\midrule
問い & どの入力を変えると、何が変わると考えたか \\
モデルと条件 & 重力源、テスト粒子近似、質量・長さ・単位、変更した入力と連動する初期値 \\
数値設定 & KDK、時間刻み、積分時間、粒子数、seed、保存間隔 \\
結果 & 同じ軸範囲の静止画と、定義を示した一つの定量指標 \\
確認 & 少なくとも一つの刻み試験、標本または測定条件への依存 \\
解釈と限界 & 分かったこと、今回の計算だけでは分からないこと \\\bottomrule
\end{tabularx}
\end{table}

例えばPPとPRの比較なら、「本書の固定形状Plummer場・同一平面のモデルでは、近点後約0.20 Gyrに、円盤2の30 kpc外側割合が逆行の場合に小さくなった」と書けます。この範囲なら、実際に測ったことと結論が対応します。一方、「逆行銀河は合体しない」「実際の銀河でこの割合の質量が失われる」という結論は、このモデルからは導けません。

設定だけでなく初期配列も保存します。乱数環境や生成手順が変わった場合に、同じseedという記録だけで完全な一致を仮定しないためです。本章の配布データには、各ケースの初期配列、出身番号、比較時刻のスナップショット、粒子割合、実行環境の版を含めています。

\section{本章のまとめ}
\label{sec:ch6-summary}
本章では、第5章で完成したKDKコードをそのまま使い、回転方向、近点距離、質量比を一つずつ変えて比較しました。入力を一つ変更すると、モデルの式から初期速度なども連動して作り直す必要があります。比較では、計算開始からの時刻だけでなく、近点からの経過時間をそろえました。

外側粒子割合は、図の見た目を補う簡単な指標です。ただし、閾値と測定時刻に依存し、非束縛率や質量損失率を直接意味しません。PPとPRの差は、今回試した閾値やseedを変えても残りました。一方、近点距離と質量比の系列では、この割合が単調に変わるとは限りませんでした。結果を予想に合わせて言い換えず、何を固定し、何を変え、何を測ったかに戻って解釈します。

本章の到達点は、完成したシミュレーションを「実行する」だけでなく、問いを立て、条件をそろえて比較し、数値的な確認とモデルの限界を含めて短い研究として説明することです。次の第\ref{chap:dynamical-friction}章では、比較実験の枠組みを保ちながら、モデルそのものへ新しい物理として動力学的摩擦を加えます。

\chapterexercises
\paragraph{演習6.1　確認：比較条件を整理する}
回転方向、近点距離、質量比の系列について、直接変える入力、固定する入力、連動して計算し直す初期値を挙げてください。質量比系列が、第2章の全質量固定の実験とは異なる理由も説明してください。

\paragraph{演習6.2　確認：内部回転だけを反転する}
円盤2の星の速度を$2\Vvec_2-\vvec_i$へ変える理由を説明してください。同じseedから生成したPRと一致するか、中心軌道と円盤1の星がPPと同じかも確認してください。この一致を一般の自己重力銀河へ適用できるでしょうか。

\paragraph{演習6.3　応用：時刻のそろえ方}
表\ref{tab:ch6-response}の近点後約0.20 Gyrと、開始から0.60 Gyrの比較を再計算してください。$q=0.1$のモデルで、$t=0.47\unit{Gyr}$は近点後何Gyrに相当するでしょうか。同じ保存番号を比べることと、同じ進化段階を比べることの違いを述べてください。

\paragraph{演習6.4　応用：指標と標本を変える}
PP--PRの円盤2について、表\ref{tab:ch6-threshold}と表\ref{tab:ch6-seed}を再計算してください。閾値の変更とseedの変更が、それぞれ何を調べる操作か説明してください。三つのseedの結果だけから「すべての銀河でこの差が成り立つ」と結論してよいでしょうか。

\paragraph{演習6.5　応用：刻みと結論の精度}
近点5 kpcのケースを0.05、0.025 Myrで計算し、$\delta_{\max}$、$\delta_{\rm rms}$、$\Delta f_k$を比べてください。$\Delta f_k=0$でも粒子位置が一致しない理由を述べてください。さらに一段細かい刻みを加える場合には、比較する時刻と初期配列をどうそろえるか説明してください。

\paragraph{演習6.6　最終課題：一つの問いを報告する}
本章の三つの系列の一つを選び、表\ref{tab:ch6-report}の順で結果をまとめてください。同じ表示範囲の図、一つの定量指標、少なくとも一つの数値的な確認を含めます。結論には、用いた近似、比較時刻、指標の定義を反映させてください。

\paragraph{演習6.7　応用：有効長さと中心間重力}
$b=3\unit{kpc}$を固定し、$D=5,8,12\unit{kpc}$で式\eqref{eq:ch6-softening-ratio}を計算してください。$D\gg b$の極限も確かめてください。この比だけで、星が受ける潮汐力の大きさや、近点系列の外側割合の順位を決められるでしょうか。また、$b$を小さくすることと、時間刻みを小さくすることの違いを説明してください。
